% !TeX root = ../main.tex
\chapter{实验评估与结果讨论}

\section{研究问题与评价指标}

实验围绕前三章提出的方法问题展开，而不是把实验设计成与方法并列的研究主题。本文考察以下问题：
\begin{description}[leftmargin=1.35cm,style=nextline]
  \item[RQ1：正确性] 压缩与解码是否保留当前输入范围内定义的结构语义，查询结果是否与独立预言机一致？
  \item[RQ2：空间] 编码中间信息复用后的列值、完整 relation、索引和共享字典分别占用多少空间？
  \item[RQ3：查询] 直接查询压缩表示及使用 GIN 时，在当前查询语义和环境下观察到怎样的延迟？
  \item[RQ4：代价] Bloom 与路径索引分别带来怎样的性能变化？现有写入测量能否评价本次测试的压缩表示？
\end{description}

相应指标包括结构差分与查询结果一致性、平均列值大小、表及 relation 大小、GIN 与共享字典字节数、查询延迟（实例中位数）以及开关 Bloom 或路径索引后的描述性差异。列值压缩率、relation 总量和查询延迟分别回答不同问题，不合并成单一“整体提升”指标。现有写入吞吐测量的表示与本次查询、存储实验不一致，故不作为本方法配置的写入评价指标。

本实验定位为原型的可行性验证和基本性能观测，不以全面复现所有 XML 压缩论文为目标，也不声称覆盖所有生产负载。

\section{实验环境、数据集与对比方案}

查询与存储实验运行于 Linux 7.3.0-9-generic x86\_64、glibc 2.44，使用启用 libxml 的 PostgreSQL 16.11（x86\_64-pc-linux-gnu，由 GCC 15.3.0 编译）；本文扩展、PostgreSQL 和比较器均以 O2 构建。三次独立运行采用 $\texttt{fsync=on}$、$\texttt{full\_page\_writes=on}$、$\texttt{synchronous\_commit=on}$、$\texttt{jit=off}$、$\texttt{max\_parallel\_workers\_per\_gather=0}$、$\texttt{default\_toast\_compression=pglz}$ 及 $\texttt{shared\_buffers=16384}$。CPU 型号、物理内存和 zlib/zstd 运行库版本未记录。

输入为按文件名排序的导航 XML 文档前缀。10,000 文档数据集的 SHA-256 为 $\texttt{124e6287a1826a20918b2c07d54b1c469c578b029f8bc34805d3457825fa57b0}$。1K、5K 和 10K 规模分别使用独立新建的 PostgreSQL 数据目录；10K 完成三次独立重复。每个查询先预热 5 次，再正式测量 10 次；报告跨实例中位数的中位数。该协议对应 warm-cache 测量，不是冷盘延迟。

查询对照包括本文压缩表示的顺序 exact-token 查询及 GIN 查询、保留原始 XML 文本并生成短叶 token 数组后建立 GIN 的工程基线、原生 XML 上语义对应的 exact-token XPath 查询，以及 zlib-9、zstd-3 和 zstd-19 压缩后逐行解压并解析 XML 的无索引扫描基线。通用压缩基线用于观测纯压缩表示上的查询处理代价，不作为索引系统对照。路径存在及路径取值另与相应消融或参照操作核验。测试 token $\texttt{neiye\_exto\_ver}$ 命中 1,816 行；缺失 token 命中 0 行。

空间对照须区分配置：原生 $\texttt{xml}$ relation 是无索引存储参考；本文压缩表示的 relation 包括表和 GIN；$\texttt{text[] + GIN}$ 基线保留原始 XML 文本，并额外物化 token 数组及其索引。该基线实现本实验的精确短 token 包含语义，不是 PostgreSQL 在原生 $\texttt{xml}$ datum 上直接建立 GIN。实验记录未包含 CPU、物理内存和部分压缩库版本，故不以当前机器信息代替缺失的实验机信息。

\section{编码正确性与存储开销分析}

\textbf{正确性结果。} 全量结构差分覆盖 10,000 个 XML 输入/输出；解码后元素顺序、容器边界和叶文本差异均为 0。每个 benchmark 实例的 44 个查询条件均与独立 Python 预言机一致。PostgreSQL 16.11 隔离实例的 14 项回归检查通过，覆盖结构语义、GIN 长 token 碰撞候选、持久化、重启和路径查询。以上结果支持已测样本与测试条件下的实现正确性，不等于对所有 XML 输入进行形式化证明。

\textbf{列值空间。} 10,000 行时本文压缩表示的平均列值为 1,818.63 B，原生 XML 列为 2,278.21 B，列级观测减少 20.17\%。此结果说明当前数据和编码配置下列值较小，不代表完整 relation 或系统总存储一定较小。

为区分压缩字节数与可查询压缩表示，实验还测量了将原始 XML 分别以 zlib-9、zstd-3、zstd-19 压缩为 $\texttt{bytea}$ 的无索引基线。10K 时三者平均列值分别为 1,637.30 B、1,645.94 B、1,524.99 B，均小于本文压缩表示的 1,818.63 B；对应无索引 relation 大小分别为 18,432,000 B、18,333,696 B、17,104,896 B。压缩基线 relation 不含查询索引，本文方法 relation 则含 GIN，因此二者完整 relation 不能作对称系统比较。压缩基线支持本实验的精确叶文本查询，但查询时需对每行解压并解析 XML，不具备本文方法的路径/内容 ID 及 GIN 键提取机制；其结果用于展示纯压缩空间参照和顺序查询代价，不用于声称本文方法压缩率优于通用压缩器。

\begin{table}[htbp]
\centering
\caption{10K 行数据的无索引压缩基线}
\label{tab:compression-baselines}
\begin{tabular}{lrrp{0.38\linewidth}}
\toprule
算法 & 平均列值（B） & relation（B） & 测试查询处理\\
\midrule
zlib-9 & 1,637.30 & 18,432,000 & 顺序扫描、逐行解压并解析 XML\\
zstd-3 & 1,645.94 & 18,333,696 & 顺序扫描、逐行解压并解析 XML\\
zstd-19 & 1,524.99 & 17,104,896 & 顺序扫描、逐行解压并解析 XML\\
\bottomrule
\end{tabular}
\end{table}

完整空间对比见表~\ref{tab:relation-space}。第一行无索引，第二、三行含索引，因此原生 XML 的 25,640,960 B 只能作为不对称配置的存储参考，不能据此判断本文含 GIN 的 relation 比含同类索引的原生 XML relation 大或小。与同测试语义的 $\texttt{text[] + GIN}$ 基线相比，本文方法完整 relation 小 50.49\%。共享字典文件另有 9,771,540 B；将一份字典全额计入后，测量持久化数据为 39,803,412 B，较该基线 relation 小 34.38\%。该核算不含数据库级及文件系统级其他开销，也未将字典在多表之间摊销；TOAST 已包含在 relation 统计中，不重复加算。

\begin{table}[htbp]
\centering
\small
\caption{不同存储与查询方案的空间占用对比}
\label{tab:relation-space}
\begin{tabular}{p{0.28\linewidth}rrrp{0.22\linewidth}}
\toprule
存储和查询方案 & 表大小（B） & GIN 大小（B） & 完整 relation（B） & 口径说明\\
\midrule
原生 PostgreSQL XML，无索引 & 25,640,960 & 0 & 25,640,960 & 原生 XML 存储参考\\
本文压缩表示 + GIN & 20,987,904 & 9,043,968 & 30,031,872 & 压缩值及 token GIN\\
原始 XML text + 生成的 text[] + GIN & 35,659,776 & 24,993,792 & 60,653,568 & 相同 exact-token 语义的工程索引基线\\
\bottomrule
\end{tabular}
\end{table}

\section{查询性能与索引效果分析}

10K warm-cache 三次独立实例的测量结果见表~\ref{tab:query-latency}。首先比较精确叶文本 token 查询：zlib/zstd 基线通过自定义函数逐行解压并解析 XML 后精确比对叶文本，均为无索引顺序扫描。其查询语义与本文方法的 exact-token 查询一致，但处理路径和索引配置不同，结果应理解为当前实现条件下的观测。

\begin{table}[htbp]
\centering
\small
\caption{10K warm-cache 查询延迟比较}
\label{tab:query-latency}
\begin{tabular}{p{0.39\linewidth}rrrr}
\toprule
查询条件 & $r_1$（ms） & $r_2$（ms） & $r_3$（ms） & 跨实例中位数的中位数（ms）\\
\midrule
本文方法顺序扫描，Bloom 开 & 72.984 & 87.450 & 73.102 & 73.102\\
本文方法顺序扫描，Bloom 关 & 85.433 & 77.111 & 100.729 & 85.433\\
本文方法 + GIN & 28.246 & 27.032 & 24.934 & 27.032\\
text[] + GIN & 27.020 & 29.911 & 24.713 & 27.020\\
zlib-9，无索引解压解析扫描 & 1,808.592 & 1,395.810 & 1,421.261 & 1,421.261\\
zstd-3，无索引解压解析扫描 & 1,455.782 & 1,329.179 & 1,429.047 & 1,429.047\\
zstd-19，无索引解压解析扫描 & 1,860.057 & 1,455.615 & 1,468.413 & 1,468.413\\
原生 XML exact-token xpath\_exists & 2,220.014 & 1,763.064 & 1,650.346 & 1,763.064\\
\bottomrule
\end{tabular}
\end{table}

本文方法与 $\texttt{text[] + GIN}$ 的中位延迟分别为 27.032 ms 和 27.020 ms，几乎相同。该结果支持“压缩值可以为 GIN 提供可用索引键，且本测试延迟与该工程基线接近”，不支持本文方法 GIN 查询更快的结论。与原生 XML exact-token XPath 的差异是本轮特定表达式、数据集和构建下的观测，不应扩展为所有 XPath 查询的普遍加速比。

通用压缩基线虽然平均列值更小，但其解压加 XML 解析的逐行扫描在本次观测中的中位数约为 1.42--1.47 s；本文方法顺序扫描中位数为 73.102 ms。该差异与编码复用减少查询阶段解压/解析工作的机制相一致，但由于原型、解析实现和压缩库执行路径不同，只能作为方法机制的初步性能观测，不能据此推导跨实现、跨数据集的普遍加速结论。

Bloom 开关下顺序查询观察到 73.102 ms 与 85.433 ms；路径索引开关下路径查询观察到 104.371 ms 与 109.618 ms。缺失 token 的 GIN 查询三轮中位数为 0.075、0.075 和 0.074 ms，表明该测试条件下字典未命中路径可快速返回；不能将其推广为所有负例或规模下的严格常数时间。路径索引开启时原始值为 104.371、107.629、104.064 ms，关闭时为 121.733、109.618、96.323 ms；对应实例中位数的中位数分别为 104.371 ms 和 109.618 ms。实验仅有三次独立实例且未进行统计显著性检验，因此这些结果仅作为描述性消融。

\section{写入开销与消融实验}

本文方法的写入过程包括 XML 解析、路径和值标识分配、共享字典访问、压缩序列化及索引维护，因而需要独立评价写入代价。目前没有与本章存储和查询实验使用同一二进制表示的写入吞吐测量。已有早期写入记录对应另一种表示，不能与本章结果合并，也不能用于判断本文所测表示相对原生 XML 或文本存储的写入性能。解析、字典访问、压缩编码与预写式日志（Write-Ahead Logging, WAL）的分项成本及并发吞吐仍待测量。

消融结果已在本章查询分析中分别报告 Bloom 和路径索引开关。在同一数据集上关闭 Bloom 后，顺序 exact-token 查询中位数从 73.102 ms 变为 85.433 ms；关闭路径索引后，路径查询中位数从 104.371 ms 变为 109.618 ms。由于重复次数有限且未做显著性检验，这些值只能说明本轮运行中的方向和量级，不能证明其他文档分布下有相同收益。共享字典是当前二进制标识契约的基础，现有实验未提供“取消共享字典但保持其他语义不变”的可比消融，因此不对该项作单独收益归因。

\section{结果分析与有效性威胁}

综合 RQ1--RQ4，当前实验支持以下条件化判断：编码阶段形成的路径和值标识能够在当前原型内被压缩表示、指定查询和 GIN 键提取共同消费；结构流、Bloom Filter 和精确 recheck 为解码及候选验证提供了约束；在所测数据和配置下，压缩列值较原生 XML 列值减少，GIN 查询延迟与 text[] + GIN 工程基线接近。上述判断描述的是方法链的可实现性及本轮空间、查询权衡，不是对所有 XML 工作负载的性能排名。

本章结果存在以下有效性威胁。第一，数据集为单一导航 XML 文档前缀，尚未覆盖不同路径深度、值分布、重复率和长 token 比例；因此空间收益不能直接推广到其他 XML 数据。第二，查询测量采用 warm-cache 协议，且实验记录缺少 CPU、物理内存及部分压缩库版本，冷盘延迟和跨机器可复现性仍需补充。第三，原生 XML、无索引通用压缩基线和含 GIN 的本文方法配置并不对称，空间结果必须按列值、relation、索引和字典分别解释。第四，当前写入测量与本章压缩表示不一致，不能据此断言本文方法的写入吞吐或更新代价。第五，原型仅实现本文定义的 XPath 导向查询操作，完整 XPath 语法、命名空间、属性轴及更多 XML 节点类型尚未在本实验中覆盖。

因此，本文支持的是“编码阶段形成的中间标识能够在当前原型内被压缩表示、指定查询和 GIN 键提取复用，并在所测条件下保持正确”的结论；不能单独证明该方法对所有数据库、任意 XML 输入或任意查询负载都更优。

\section{本章小结}

本章按照正确性、空间、查询和写入成本评价本文方法。正确性检查在已测文档和查询条件上通过；本文压缩列值小于原生 XML，但大于无索引 zlib/zstd 压缩基线，后者 relation 不含索引，不能与本文含 GIN 的配置作对称比较。与同语义 text[] + GIN 基线相比，GIN 延迟基本相当，计入一份字典后的测量持久化成本仍较低。当前尚无与本章表示匹配的写入吞吐证据。所有定量结论均限制于对应数据、环境和测试协议。
